\title{Controlling Text-to-Image Diffusion by Orthogonal Finetuning}
\textbf{Convergence}. To assess how rapidly each method learns, we investigate the convergence rates of ControlNet, T2I-Adapter, LoRA, and COFT on the L2F task, employing both quantitative and qualitative measures. For the quantitative aspect, we calculate the mean $\ell_2$ distance between the ground truth face landmarks specified by control and those predicted by the models. Figure~\ref{fig:face_convergence} presents the trajectory of face landmark errors over different training epochs. These results reveal that COFT and OFT reach low error values notably sooner than the alternatives. For instance, LoRA needs 20 epochs to attain the same performance level that OFT achieves by epoch 8. It is important to highlight that both OFT and COFT involve a comparable parameter count to LoRA—substantially less than ControlNet—yet display much faster convergence than the existing baselines. This rapid convergence characteristic of OFT is corroborated by the evidence in Figure~\ref{fig:teaser}, where the example on the right demonstrates OFT's superior data efficiency relative to ControlNet and LoRA; OFT maintains strong convergence even with just 5\% of the ADE20K data. Regarding qualitative analysis, our comparison emphasizes OFT, COFT, and ControlNet, as ControlNet produces landmark errors closest to ours. Figure~\ref{fig:face_convergence_qual} illustrates that OFT and COFT both exhibit stable convergence, with generated face poses progressively aligning with the prescribed landmarks. Conversely, ControlNet displays a delayed onset of controllability—becoming apparent only after epoch 8—a pattern aligning with the abrupt convergence behavior described in \cite{zhang2023adding}. This degree of stability makes OFT especially practical, as its training dynamics are more gradual and reliable, facilitating easier selection of appropriate hyperparameters.

\section{Introduction}

\section{Related Work}

The rationale behind employing orthogonal transformation for finetuning neurons is partly inspired by the mechanics of the 2D Fourier transform, which decomposes an image into its magnitude and phase components. Notably, the phase spectrum encapsulates angular relationships between the input and the basis, and is responsible for retaining much of the image’s semantic content. As a case in point, combining the phase spectrum of a given image with an arbitrary magnitude spectrum can reconstruct an image that still conveys the original semantics. This observation implies that altering neuron orientations predominantly modulates the semantics of the generated images—central to tasks in subject-driven and controllable generation. Nonetheless, allowing unrestricted changes in neuron directions risks undermining the generative capability established during pretraining. To address this, we introduce a constraint that preserves the angles between every pair of neurons, motivated by the premise that these inter-neuron angles play a vital role in capturing the neural network’s encoded knowledge. Based on this, we propose a layer-wise shared orthogonal transformation that ensures the hyperspherical energy remains constant across finetuning.

    \item We introduce Orthogonal Finetuning, a new approach for adapting text-to-image diffusion models to enhance controllability. Additionally, we present a constrained variant that restricts the angular shift relative to the original pretrained model in order to promote greater training stability.
    \item In contrast to standard finetuning approaches, OFT preserves the hyperspherical energy during adaptation—something we rigorously demonstrate and empirically validate as a critical indicator of retaining the semantic generative capabilities of the base model.
    \item Our method is evaluated on both subject-driven generation and controllable generation tasks. Through extensive empirical analysis, we report marked improvements over previous techniques regarding output quality, training convergence rate, and finetuning robustness. Furthermore, OFT demonstrates superior data efficiency, achieving strong performance using considerably fewer training examples and epochs.
    \item For the controllable generation setting, we propose a novel control consistency metric to quantify controllability. The core procedure estimates the control signal inferred from generated outputs, subsequently comparing it to the original control input. Our results confirm that OFT offers enhanced control over the generative process.
\end{itemize}

\textbf{Why not minimize hyperspherical energy}. Unlike the approach in \cite{liu2021orthogonal}, our goal is not to minimize hyperspherical energy. In classification settings, redundancy-free neuron arrangements are preferable, and achieving minimum hyperspherical energy implies that all neurons are evenly distributed over the hypersphere. However, such uniform spacing is not a suitable finetuning objective because it risks undermining the knowledge obtained through pretraining.

\begin{itemize}[leftmargin=*,nosep]

\section{Intriguing Insights and Discussions}

Inspired by the pivotal role that angular relationships among neurons play in shaping visual semantics, we introduce a straightforward yet powerful finetuning approach—orthogonal finetuning—for exerting control over text-to-image diffusion models. Our focus is on two primary text-to-image use cases: subject-driven generation and controllable generation. When benchmarked against previous methods, OFT achieves improved controllability and greater finetuning robustness while requiring fewer finetuning parameters. Notably, OFT imposes no additional computational cost during inference, thus making it highly practical for deployment.

\subsection{Efficient Orthogonal Parameterization}\label{sect:eff_ortho}

\subsection{Constrained Orthogonal Finetuning}

Typically, standard finetuning updates the parameters of a neural network—either all weights or select layers—by applying gradient descent with a modest learning rate. This conservative update strategy implicitly ensures that the finetuned weights do not stray far from their pretrained counterparts, thereby minimizing $\thickmuskip=2mu \medmuskip=2mu \| \bm{M} - \bm{M}^0\|$, where $\bm{M}$ and $\bm{M}^0$ denote the adapted and original model weights, respectively. While this indirect regularization is intuitively appealing, it may not sufficiently constrain the modifications when adapting very large models. To further restrict updates, LoRA explicitly adds a low-rank constraint to the difference between updated and initial weights, formalized as $\thickmuskip=2mu \medmuskip=2mu \text{rank}(\bm{M} - \bm{M}^0)=r'$, for some small rank $r'$. In contrast, OFT enforces a different restriction, preserving the pairwise relationships among neurons through their hyperspherical energy: $\thickmuskip=2mu \medmuskip=2mu \| \text{HE}(\bm{M})-\text{HE}(\bm{M}^0) \|=0$, with $\text{HE}(\cdot)$ denoting the hyperspherical energy measure for a matrix. To illustrate, consider a dense layer with weight matrix $\thickmuskip=2mu \medmuskip=2mu \bm{W}=\{\bm{w}_1,\cdots,\bm{w}_n\}\in\mathbb{R}^{d\times n}$, where each column vector $\bm{w}_i\in\mathbb{R}^{d}$ is the $i$-th neuron's parameters and $\bm{W}^0$ represents the pretrained weights. The output vector $\thickmuskip=2mu \medmuskip=2mu \bm{z}\in\mathbb{R}^n$ is calculated as $\thickmuskip=2mu \medmuskip=2mu \bm{z}=\bm{W}^\top\bm{x}$ given the input $\thickmuskip=2mu \medmuskip=2mu \bm{x}\in\mathbb{R}^d$. The goal of OFT is to minimize the difference in hyperspherical energy between the updated and the initial weights:

\textbf{Connection to LoRA}. LoRA maintains the integrity of pretrained weights by introducing a low-rank matrix, ensuring that the information stored in the original parameters remains mostly intact. By contrast, OFT imposes an orthogonal (and thus full-rank) transformation on the pretrained weight matrix, which safeguards against substantial degradation of pretraining information. The forward propagation for OFT can be expressed as $\thickmuskip=2mu \medmuskip=2mu \bm{z}=(\bm{R}\bm{W}^0)^\top\bm{x}=(\bm{W}^0+(\bm{R}-\bm{I})\bm{W}^0)^\top\bm{x}$, where $\thickmuskip=2mu \medmuskip=2mu (\bm{R}-\bm{I})\bm{W}^0$ serves a similar role to the low-rank modification used in LoRA. Typically, since $\bm{W}^0$ is full-rank, OFT effectively conducts a low-rank update whenever $\thickmuskip=2mu \medmuskip=2mu \bm{R}-\bm{I}$ has low rank. Comparable to LoRA’s tunable rank $r'$, OFT features a block diagonal parameter $r$, allowing a reduction in the trainable parameter count. Notably, the two methods embody fundamentally different strategies for parameter efficiency: LoRA leverages low-rankness, whereas OFT capitalizes on the sparsity brought by block-diagonal orthogonality constraints.

\begin{equation}\label{eq:oft_general}
    \footnotesize
    \bm{z}=\bm{W}^\top\bm{x}=(\bm{R}\cdot\bm{W}^0)^\top\bm{x},~~~\text{s.t.}~\bm{R}^\top\bm{R}=\bm{R}\bm{R}^\top=\bm{I}
\end{equation}

However, OFT presents several intriguing unresolved questions. Firstly, orthogonality within OFT is maintained using Cayley parametrization, a technique that necessitates matrix inversion, slightly restricting the approach’s scalability. We partially mitigate this constraint by leveraging block diagonal parametrization; nevertheless, discovering efficient, differentiable methods for matrix inversion remains an open challenge. Secondly, OFT is particularly promising for compositionality: orthogonal matrices generated from different OFT finetuning processes can be multiplied to yield another orthogonal matrix. A worthwhile direction for future exploration is whether such a collection of orthogonal matrices can simultaneously preserve the distinct knowledge from all the respective downstream tasks. Lastly, OFT's parameter efficiency is intimately tied to its block diagonal parameterization, which can introduce structural biases and reduce modeling flexibility. Enhancing parameter efficiency with methods that are more adaptive and introduce less bias is an open research avenue of significant interest.

    \item \textbf{Subject-driven generation}~\cite{ruiz2023dreambooth}: Starting from a handful of images featuring a specific subject, the objective is to synthesize new depictions of that subject in novel settings guided by a text prompt.
    \item \textbf{Controllable generation}~\cite{zhang2023adding,mou2023t2i}: By providing an auxiliary control input (\emph{e.g.}, edge maps, segmentation masks), the goal is to guide the model to create images adhering to the given control while also following a text prompt.
\end{itemize}

\textbf{OFT is architecture-independent}. In theory, OFT is applicable to all neural network models. For example, within Transformers, LoRA customarily targets the attention weights~\cite{hulora2022}. For an equitable comparison to LoRA, our experiments also restrict OFT application to the finetuning of attention weights. Beyond fully connected layers, OFT offers advantages for convolutional layer finetuning: the block-diagonal form of $\bm{R}$ acquires specific interpretability when dealing with convolutional layers, a feature not present with LoRA. In particular, matching the number of blocks to the input channels ensures each block modulates one neuron channel individually, much like the principles behind depth-wise convolution~\cite{chollet2017xception}. If, instead, every block in $\bm{R}$ is shared, a unified orthogonal transformation operates over all channels. Further exploratory work on applying OFT to convolutional layer finetuning can be found in Appendix~\ref{app:convolution}.

\textbf{Quantitative comparison}. To assess the effectiveness of controllable generation, we propose a metric for control consistency. This approach involves extracting the control signal from the generated output and measuring its correspondence to the original control input. For the C2I task, we report IoU and F1 score; for S2I, mean IoU, mean accuracy, and overall accuracy; and for L2F, the average $\ell_2$ distance between the input control landmarks and those predicted in the generated images. Detailed explanations of the employed consistency metrics are available in Appendix~\ref{app:settings}. Each method is evaluated using its optimal hyperparameters. According to Table~\ref{table:control_gen}, OFT and COFT consistently outperform competing approaches in terms of strength and precision of control. Notably, adapter-based methods (\emph{e.g.}, T2I-Adapter, ControlNet) exhibit slower convergence and achieve inferior results at convergence. While LoRA surpasses ControlNet on S2I, it underperforms on C2I and L2F. Generally, even with extensive hyperparameter optimization, LoRA demonstrates limited potential. In contrast, we find the OFT framework’s performance continues to improve as more parameters are made trainable, indicating that it has not yet reached its limit. We highlight our quantitative protocol for controllable generation as an important contribution, enabling precise measurement of control quality in finetuned models (within the bounds set by the segmentation/detection models employed).

While classical orthogonalization procedures like the Gram-Schmidt process are differentiable, they generally incur significant computational cost in real-world settings~\cite{liu2021orthogonal}. To improve efficiency, we utilize the Cayley transform for generating the orthogonal matrix. More concretely, the orthogonal matrix is parameterized as $\thickmuskip=2mu \medmuskip=2mu \bm{R}=(\bm{I}+\bm{Q})(I-\bm{Q})^{-1}$, where $\bm{Q}$ is a skew-symmetric matrix such that $\thickmuskip=2mu \medmuskip=2mu \bm{Q}=-\bm{Q}^\top$. This method achieves computational gains, though with a minor trade-off: the Cayley parameterization is limited to producing orthogonal matrices whose determinant equals $1$, meaning $\bm{R}$ is confined to the special orthogonal group. Empirically, we observe this restriction does not hinder performance. Even with the Cayley-based parameterization, representing a full $\bm{R}$ remains inefficient in terms of parameters when $d$ is large. To overcome this, we introduce a block-diagonal structure for $\bm{R}$, comprising $r$ blocks. This leads to the formulation:

\subsection{Why Does Orthogonal Transformation Make Sense?}

\textbf{Controllable generation}. Image-to-image translation represents a paradigm of controllable generation. Traditionally, this has been addressed via conditional generative adversarial networks~\cite{isola2017image,zhu2017toward,wang2018high,park2019semantic,choi2018stargan}. More recently, diffusion models have been applied in this context~\cite{saharia2022palette,voynov2022sketch,wang2022pretraining}. Among notable advances, ControlNet~\cite{zhang2023adding} enhances controllability of pretrained diffusion models by finetuning and incorporating external control signals, achieving remarkable generation control. Similarly, T2I-Adapter~\cite{mou2023t2i} fine-tunes pretrained diffusion models to offer stronger guidance over generated outputs. Adopting the framework in \cite{zhang2023adding,mou2023t2i}, our approach utilizes OFT for finetuning pretrained diffusion models. This yields consistently superior controllability while requiring less training data and fewer finetuned parameters. Notably, OFT also incurs no additional computational cost at test-time inference.

\textbf{Qualitative comparison}. To provide a clearer picture of the advantages offered by OFT, we present a selection of randomly chosen instances of subject-driven generation in Figure~\ref{fig:dreambooth_final}. To ensure a just evaluation, all samples are derived from a single finetuned model per method, with no independent prompt tuning for each outcome. For every approach, the model demonstrating the highest validation CLIP scores is employed. Inspection of the outcomes in Figure~\ref{fig:dreambooth_final} reveals that both OFT and COFT excel at maintaining the semantic features of the target subject, whereas LoRA frequently fails to retain subject identity (for instance, LoRA entirely omits the subject identity in the bowl scenario). Furthermore, OFT and COFT allow for more faithful control via text instructions, whereas DreamBooth—despite strong subject identity retention—regularly does not align the generated content with the specified prompt (e.g., in the first row featuring the bowl). This qualitative analysis highlights OFT's ability to simultaneously deliver superior controllability and faithful subject preservation. Additionally, OFT's performance is robust to the number of optimization iterations; it maintains effectiveness even with extensive iteration counts, a property not shared by either DreamBooth or LoRA. Further illustrative results can be found in Appendix~\ref{app:more_qualitative}. We also provide a human evaluation in Appendix~\ref{app:human} to further corroborate the strengths of OFT.

At their core, both tasks address the challenge of how to adapt text-to-image diffusion models to new objectives through finetuning, while maintaining the generative quality established during pretraining. The requirements for a successful finetuning approach can be characterized as follows: (1) \emph{training efficiency}, which involves minimizing the number of parameters that require updating as well as reducing the training duration, and (2) \emph{generalizability preservation}, ensuring the finetuned model retains its ability to produce high-quality and varied outputs. Common approaches include directly adjusting the model's parameters with a small learning rate (\emph{e.g.}, \cite{ruiz2023dreambooth}), or augmenting the architecture with a lightweight module containing re-parameterized neuron weights (\emph{e.g.}, \cite{hulora2022,zhang2023adding}). While these methods are straightforward, neither can unequivocally ensure that the generative capacity achieved during pretraining remains intact. Moreover, there is a notable absence of systematic guidelines for designing effective finetuning schemes and selecting appropriate hyperparameters, such as the number of training epochs. One significant obstacle is the lack of an established metric to assess how much of the pretrained generative function has been preserved. Current finetuning techniques often operate under the implicit hypothesis that minimizing the Euclidean distance between the pretrained and finetuned models serves as an adequate proxy for retention of generative abilities, leading to the use of small learning rates. However, while this rationale sometimes holds, we contend that Euclidean proximity by itself fails to thoroughly represent the preservation of semantic content. Consequently, introducing a more structurally aware metric to describe differences between the pretrained and finetuned models would substantially enhance both the stability and efficacy of finetuning with respect to pretraining performance.

\begin{itemize}[leftmargin=*,nosep]

\textbf{Quantitative comparison}. Adopting the methodology of \cite{ruiz2023dreambooth}, we quantitatively benchmark subject fidelity (DINO~\cite{caron2021emerging}, CLIP-I~\cite{radford2021learning}), text prompt fidelity (CLIP-T~\cite{radford2021learning}), and sample diversity (LPIPS~\cite{zhang2018unreasonable}). CLIP-I measures the mean pairwise cosine similarity between CLIP embeddings of generated and reference images; DINO adopts the same approach using ViT S/16 DINO features instead. For CLIP-T, we compute the average cosine similarity between CLIP embeddings of the textual prompt and the generated images. We further assess sample diversity using the mean LPIPS cosine similarity between generations of the same subject given an identical prompt. The outcomes, reported in Table~\ref{table:dreambooth_final}, indicate that COFT and OFT considerably surpass DreamBooth and LoRA on DINO and CLIP-I metrics, while they either match or modestly exceed these baselines with respect to prompt fidelity and diversity. To reduce stochasticity, each method undergoes finetuning 30 times from the same pretrained model with distinct random seeds. Collectively, the evidence suggests our OFT framework achieves not only superior convergence and stability, but also consistently outperforms baselines in final quantitative metrics.

\subsection{General Framework}

\begin{equation}\label{eq:oft_principle}
\footnotesize
\min_{\bm{W}} \left\| \text{HE}(\bm{W})-\text{HE}(\bm{W}^0)\right\|~~\Leftrightarrow~~\min_{\bm{W}} \bigg{\|} \sum_{i\neq j} \|\hat{\bm{w}}_i-\hat{\bm{w}}_j\|^{-1}-\sum_{i\neq j} \|\hat{\bm{w}}^0_i-\hat{\bm{w}}^0_j\|^{-1}\bigg{\|}
\end{equation}

With respect to the second question, a straightforward strategy to adjust neuron orientations is to apply a collective rotation or reflection (within the same layer) to all neurons, naturally corresponding to an orthogonal transformation. Although alternative angular manipulations could introduce greater flexibility by assigning unique rotation angles to different neurons, our experiments indicate that orthogonal transformations strike an optimal balance between expressiveness and structural coherence. Additionally, \cite{liu2021orthogonal} provides evidence that orthogonal transformations possess sufficient expressive power for training neural networks. To substantiate this claim, we conduct an experiment measuring the regularizing effects induced by enforcing orthogonality, employing the ControlNet~\cite{zhang2023adding} framework for controllable generation; the results are shown in Figure~\ref{fig:ortho}. The findings show that the canonical OFT method, when equipped with the orthogonality constraint, produces stable and precise control by the end of training (epoch 20), while removing this constraint results in an inability to produce meaningful outputs or control. These results affirm the necessity of maintaining orthogonality in OFT.

\textbf{Qualitative comparison}. We further conduct a qualitative assessment of OFT and COFT against leading approaches such as ControlNet, T2I-Adapter, and LoRA. As depicted in Figure~\ref{fig:control_final}, random samples illustrate that both OFT and COFT produce images with superior realism and fidelity, in addition to enabling precise control. Within the S2I task, LoRA fails to generate outputs aligned with the provided segmentation maps, whereas ControlNet, OFT, and COFT are capable of preserving the intended structures in their results. Compared to ControlNet, our OFT and COFT methods generate images exhibiting finer details and more plausible geometrical coherence, while utilizing significantly fewer parameters. For the C2I task, OFT and COFT convincingly create images that are semantically consistent with the coarse Canny edges, whereas T2I-Adapter and LoRA demonstrate inferior performance. Regarding the L2F task, our method most effectively maintains accurate pose information in generated faces, even when presented with difficult pose variations. Across all three control tasks, OFT and COFT consistently deliver images of higher quality compared to state-of-the-art competitors, underscoring the strengths of our proposed OFT framework for controllable image generation. For an expanded qualitative comparison, we offer additional examples for each control task in Appendix~\ref{app:control_exp}, and further showcase the versatility of OFT in handling other control settings—such as dense pose to human body, sketch to image, and depth to image—in Appendix~\ref{app:more_control_task}.

\section{Experiments and Results}

\textbf{General experimental setup}. Our experiments utilize Stable Diffusion v1.5~\cite{rombach2022high} as the pretrained model for text-to-image generation. For unbiased assessment, images generated by different techniques are sampled at random. The protocol for subject-driven generation aligns with DreamBooth~\cite{ruiz2023dreambooth}, while for controllable generation, the methodologies from ControlNet~\cite{zhang2023adding} and T2I-Adapter~\cite{mou2023t2i} are followed. For fair comparison with LoRA, OFT or COFT are applied exclusively to layers where LoRA is also implemented. Additional experimental details and results are presented in Appendix~\ref{app:settings}.

\textbf{Subject-driven generation}. Efforts to prevent unwanted alteration of the subject include methods that condition on user-provided masks~\cite{avrahami2022blended,nichol2022glide}. Inversion-based strategies~\cite{choi2021ilvr,dhariwal2021diffusion,ramesh2022hierarchical,gal2022image} are utilized to modify scene context while maintaining subject integrity. The method in~\cite{hertz2022prompt} offers editing capabilities at both local and global scales without relying on masks, though these approaches often struggle with preserving fine identity attributes. Pivotal Tuning~\cite{roich2022pivotal} fine-tunes the generator close to an inverted latent representation with additional constraints for identity maintenance. Similarly,~\cite{nitzan2022mystyle} develops individualized generative face priors using datasets comprised of a person's facial images. The approach in~\cite{casanova2021instance} produces diverse instance variations but sometimes neglects instance-specific characteristics. DreamBooth~\cite{ruiz2023dreambooth} achieves personalized adaptation by introducing a custom token and a small number of subject images, applying both reconstruction and class-preservation objectives during finetuning. Building upon DreamBooth, OFT replaces conventional finetuning with small-step updates by leveraging orthogonal transformations for adaptation.

\end{document}

\textbf{Why OFT converges faster}. As demonstrated in Figure~\ref{fig:toy_ae}, the most effective way to alter semantic content is by adjusting neuron directions—precisely the mechanism enabled by OFT. Furthermore, OFT reframes finetuning as modifying neuron representations over a smooth hyperspherical manifold, resulting in a more favorable optimization landscape. Empirical results from \cite{liu2021orthogonal} further support this advantage.

Here, $\text{O}(d)$ represents the group of orthogonal transformations in $d$ dimensions. The matrices $\bm{R}\in\mathbb{R}^{d\times d}$ and, for each $i$, $\bm{R}_i\in\mathbb{R}^{d/r\times d/r}$ are assumed to be orthogonal. When $r=1$, the block-diagonal structure simplifies to a fully unconstrained orthogonal matrix. The total number of parameters for a general $d\times d$ orthogonal matrix is $d(d-1)/2$, giving a parameter complexity of $\mathcal{O}(d^2)$. For an orthogonal matrix with an $r$-block diagonal structure, there are $d(d/r-1)/2$ parameters, or a complexity of $\mathcal{O}(d^2/r)$. Optionally, parameter sharing among blocks—i.e., enforcing $\bm{R}_i = \bm{R}_j$ for all $i\neq j$—can further lower the parameter count to $\mathcal{O}(d^2/r^2)$. It is important to note that despite these parameter reductions, the global matrix $\bm{R}$ still remains orthogonal, thereby guaranteeing preservation of hyperspherical energy.

In this formulation, $\bm{R}$ is not only required to be orthogonal but must also lie within an $\epsilon$-ball of the identity matrix. Orthogonality can be imposed via Cayley parameterization as discussed in Section~\ref{sect:eff_ortho}, but ensuring the $\epsilon$-deviation property under this parameterization remains challenging. To shed light on the Cayley transformation, we may leverage the Neumann series, expressing $\bm{R}=(\bm{I}+\bm{Q})(I-\bm{Q})^{-1}$, and approximate it as $\bm{R}\approx\bm{I}+2\bm{Q}+\mathcal{O}(\bm{Q}^2)$, under the presumption that the Neumann

\textbf{Balancing flexibility and regularity in finetuning}. Our investigation reveals an intrinsic balance between flexibility and regularity when finetuning models. While standard finetuning offers maximum adaptability, it is often plagued by instability and can precipitate model collapse. In contrast, the simplicity of OFT achieves a favorable compromise between these factors by maintaining the pairwise angles between neurons. Moreover, the block-diagonal approach can be interpreted as imposing an even stronger constraint on the orthogonal transformation matrix.

\begin{equation}
    \footnotesize
    \bm{R}=\text{diag}(\bm{R}_1,\bm{R}_2,\cdots,\bm{R}_r)=\begin{bmatrix}
    \bm{R}_1\in \text{O}(\frac{d}{r}) &  &\\
    &\ddots & \\
    & & \bm{R}_r\in \text{O}(\frac{d}{r})
    \end{bmatrix}\in \text{O}(d)
\end{equation}

Drawing from empirical findings that hyperspherical similarity effectively captures semantic relationships~\cite{liu2017deep,liu2018decoupled,chen2020angular}, our approach utilizes hyperspherical energy~\cite{liu2018learning} to represent the pairwise interactions among neurons. Specifically, hyperspherical energy refers to the aggregate of hyperspherical similarity (\emph{e.g.}, cosine similarity) between every pair of neurons within a layer, thereby reflecting how evenly distributed the neurons are on the unit hypersphere~\cite{liu2021learning}. We propose that optimal finetuning should result in minimal changes to the hyperspherical energy relative to the original pretrained model. An intuitive strategy might be to introduce a regularization term that maintains constant hyperspherical energy throughout finetuning; however, this does not reliably ensure that the energy deviation remains small. Instead, we exploit a key invariance property: pairwise hyperspherical similarities remain unchanged under a common orthogonal transformation applied across all neurons. Leveraging this invariance, we introduce Orthogonal Finetuning~(OFT), a method designed to adapt large text-to-image diffusion models for specific downstream applications without altering their hyperspherical energy. The main concept involves learning an orthogonal transformation shared across layers, thus preserving the pairwise angular relationships between neurons. This method can also be interpreted as modifying the standard coordinate system of the neurons within a layer. Moreover, considering the observation that a reduced Euclidean distance between finetuned and pretrained models corresponds to better retention of pretraining benefits, we further develop a variant called Constrained Orthogonal Finetuning~(COFT), which limits the finetuned model to remain on a hypersphere of fixed radius centered at the positions of the pretrained neurons.

\textbf{Settings}. The baseline approaches in our evaluation include ControlNet~\cite{zhang2023adding}, T2I-Adapter~\cite{mou2023t2i}, and LoRA~\cite{hulora2022}. In the main text, we explore three difficult controllable generation benchmarks: Canny edge to image~(C2I) on COCO~\cite{lin2014microsoft}, segmentation map to image~(S2I) on ADE20K~\cite{zhou2017scene}, and landmark to face~(L2F) on CelebA-HQ~\cite{karras2018progressive,xia2021tedigan}. All compared methods are utilized to finetune Stable Diffusion~(SD) v1.5 for 20 training epochs on each dataset. Additional experimental findings can be located in Appendix~\ref{app:more_qualitative},~\ref{app:more_control_task},~\ref{app:failure}.

\subsection{Subject-driven Generation}

In terms of parameter efficiency, we contrast OFT and LoRA. A LoRA adaptation with low rank $r'$ contains $r'(d+n)$ trainable variables. Assuming that both $r$ (in OFT) and $r'$ (in LoRA) scale with the neuron dimensionality, such that $r = r' = \alpha d$ for a constant $0 < \alpha \leq 1$, LoRA exhibits a parameter complexity of $\mathcal{O}(d^2 + dn)$, whereas for OFT the complexity reduces to $\mathcal{O}(d)$. This complexity distinction can be illustrated concretely: for a weight matrix of dimension $128\times 128$, LoRA with $r'=8$ possesses $2,048$ trainable variables, while OFT using $r=8$ (without any block sharing) requires only $960$ parameters.

Here, $\bm{W}$ is the weight matrix updated via OFT-finetuning, and $\bm{I}$ represents the identity matrix. Figure~\ref{fig:block} provides a depiction of OFT. Analogous to the zero initialization strategy used in LoRA, it is essential for OFT to ensure that the fine-tuning process begins precisely from the pretrained matrix $\bm{W}^0$. To this end, we set the initial value of the orthogonal matrix $\bm{R}$ to the identity matrix, thereby ensuring that the initial parameters match those of the pretrained model. The orthogonality of $\bm{R}$ can be preserved through differentiable orthogonalization techniques introduced in \cite{liu2021orthogonal,lezcano2019cheap}. In what follows, we examine methods to enforce orthogonality efficiently.

\textbf{Model finetuning}. Adapting large pretrained models for downstream applications has become a widespread practice~\cite{devlin2018bert,brown2020language,he2022masked}. A subset of these techniques, known as adaptation methods (\emph{e.g.}, \cite{hulora2022,houlsby2019parameter,pfeiffer2020adapterfusion}), is particularly prominent in natural language processing. OFT is most closely related to LoRA~\cite{hulora2022}, which operates under the assumption that finetuning can be effectively achieved via low-rank additive updates to weights. By contrast, OFT employs a layer-shared orthogonal transformation for multiplicative weight updates at the neuron level, ensuring the preservation of pair-wise angular relationships between neurons within the same layer and resulting in enhanced stability.

The expressive capacity of the basic OFT can be further narrowed by enforcing that the fine-tuned model stays within a limited range of the pretrained one. To that end, COFT modifies the model’s forward pass as follows:

To motivate the use of orthogonal transformation for finetuning text-to-image diffusion models, we break down the rationale into two aspects: (1) the benefits of adjusting the angular direction (i.e., the orientation) of neuron weights, and (2) the appropriateness of using orthogonal transformation for this angular adjustment.

\section{Concluding Remarks and Open Problems}

\section{Orthogonal Finetuning}

Here, $\thickmuskip=2mu \medmuskip=2mu \hat{\bm{w}}_i:=\bm{w}_i/\|\bm{w}_i\|$ is the normalized vector of the $i$-th neuron, and hyperspherical energy is defined for the layer as $\thickmuskip=2mu \medmuskip=2mu \text{HE}(\bm{W}):= \sum_{i\neq j} \|\hat{\bm{w}}_i-\hat{\bm{w}}_j\|^{-1}$. The minimum attainable value for Eq.~\eqref{eq:oft_principle} is evidently zero. This is achievable if and only if the post-finetuning weights differ from the initial weights solely by a rotation or reflection, that is, when $\thickmuskip=2mu \medmuskip=2mu \bm{W}=\bm{R}\bm{W}^0$ for some orthogonal matrix $\thickmuskip=2mu \medmuskip=2mu \bm{R}\in\mathbb{R}^{d\times d}$ (where $\det(\bm{R})=1$ or $-1$ indicates a rotation or a reflection, respectively). OFT leverages this insight by restricting adaptation to the learning of a single orthogonal transformation per layer, preserving the geometric relations among neurons. Formally

Recently developed text-to-image diffusion models~\cite{saharia2022photorealistic,ramesh2022hierarchical,rombach2022high} have set new standards in generating high-resolution images based on textual input. However, reliance on text prompts alone often results in ambiguity and lacks the granularity required for precise image control. To overcome these challenges, this work focuses on two categories of text-to-image generation problems:

\begin{equation}\label{eq:coft_general}
    \footnotesize
    \bm{z}=\bm{W}^\top\bm{x}=(\bm{R}\cdot\bm{W}^0)^\top\bm{x},~~~\text{s.t.}~\bm{R}^\top\bm{R}=\bm{R}\bm{R}^\top=\bm{I},~\norm{\bm{R}-\bm{I}}\leq \epsilon
\end{equation}

\textbf{Text-to-image diffusion models}. Substantial advances~\cite{saharia2022photorealistic,ramesh2022hierarchical,rombach2022high,gu2022vector,nichol2022glide} in text-to-image synthesis have emerged, primarily propelled by innovations in diffusion-based generative modeling~\cite{ho2020denoising,song2021score,song2021denoising,dhariwal2021diffusion} and multimodal representation learning~\cite{su2020vlbert,lu2019vilbert,radford2021learning,wang2021simvlm,li2022blip,li2023blip,alayrac2022flamingo,schuhmann2022laion}. In the pixel domain, models such as GLIDE~\cite{nichol2022glide} and Imagen~\cite{saharia2022photorealistic} have demonstrated success, with GLIDE performing joint training of the text encoder and diffusion prior, whereas Imagen employs a fixed, pretrained text encoder throughout. Meanwhile, Stable Diffusion~\cite{rombach2022high} and DALL-E2~\cite{ramesh2022hierarchical} operate within a latent space framework—Stable Diffusion utilizes VQ-GAN~\cite{esser2021taming} to establish a visual codebook as the latent space, while DALL-E2 leverages CLIP~\cite{radford2021learning} to align image and text in a shared latent space. Besides diffusion models, other paradigms such as generative adversarial networks~\cite{reed2016generative,zhang2017stackgan,xu2018attngan,li2019controllable} and autoregressive techniques~\cite{ramesh2021zero,ding2021cogview,wu2022nuwa,yuscaling} have also contributed to the field of text-conditional image synthesis. OFT is agnostic to the specific generative architecture and is compatible with any text-to-image diffusion system.

\subsection{Re-scaled Orthogonal Finetuning}

\subsection{Controllable Generation}

To address the first question, we take guidance from empirical findings in \cite{liu2018decoupled,chen2020angular}, which indicate that differences in feature angles effectively capture the semantic disparity. SphereNet~\cite{liu2017deep} demonstrates that normalizing all neurons of a neural network to lie on a unit hypersphere preserves model capacity and can even enhance generalization, suggesting that neuron orientation encodes the critical information within the data. To further clarify the significance of neuron angles, we design a toy experiment as illustrated in Figure~\ref{fig:toy_ae}, wherein a conventional convolutional autoencoder is trained using flower image data. During training, the standard inner product is applied to create the feature map (here, $z$ denotes the element-wise result of convolving kernel $\bm{w}$ with the input $\bm{x}$). For evaluation, we consider three approaches for generating the feature map: (a) the inner product (matching the training phase), (b) relying solely on magnitude, and (c) utilizing only angular information. The reconstructions depicted in Figure~\ref{fig:toy_ae} reveal that the angular component nearly reconstructs the original images, whereas the magnitude alone does not capture useful structure. It is important to note that cosine activation (c) is not used during training, where only the inner product guides learning. These findings suggest that neuron directions encode most of the semantic content in the images, and thus altering neuron directions could be the most direct means to manipulate semantic features.

Additionally, our approach draws upon techniques from orthogonal over-parameterized training~\cite{liu2021orthogonal}, where classification neural networks are initialized randomly and then trained through orthogonal transformations. This technique is grounded in the property that randomly initialized networks exhibit low hyperspherical energy on average, and the objective in \cite{liu2021orthogonal} is to maintain low hyperspherical energy during training, which empirical studies show is beneficial for generalization in classification~\cite{liu2018learning,lin2020regularizing}. While orthogonal transformations have proven sufficiently expressive for learning generalizable classification models~\cite{liu2021orthogonal}, our focus diverges as we address finetuning in text-to-image diffusion models, emphasizing enhanced controllability and improved generative performance in downstream tasks. We point out two primary distinctions between our OFT approach and the method in \cite{liu2021orthogonal}. Firstly, unlike \cite{liu2021orthogonal}, which is aimed at reducing hyperspherical energy, OFT explicitly preserves the pretrained model’s hyperspherical energy to protect its intrinsic structure during finetuning—a crucial aspect for diffusion models, where excessive minimization could compromise core semantic representations. Secondly, OFT is formulated to minimize departure from the pretrained parameters, hence introducing a constrained version. Conversely, the methodology in \cite{liu2021orthogonal} is unconstrained in this respect. Finetuning ultimately involves a delicate balance between model adaptability and preservation of existing capabilities, and we argue that OFT provides an effective solution for navigating this balance. Our primary contributions are summarized as follows:

\textbf{Setup}. DreamBooth~\cite{ruiz2023dreambooth} and LoRA~\cite{hulora2022} are adopted as reference baselines. The same loss objective as DreamBooth is used across all approaches. Both DreamBooth and LoRA are implemented following their respective original protocols, utilizing the optimal reported hyperparameters. Further experimental outcomes are detailed in Appendix~\ref{app:settings},\ref{app:coft_oft},\ref{app:more_qualitative},\ref{app:failure}.

We introduce a straightforward modification of OFT by introducing a learnable scaling factor for each neuron. This design stems from the observation that neuron rescaling does not affect hyperspherical energy since normalization removes the effect of magnitude. Specifically, the forward computation is given by $\thickmuskip=2mu \medmuskip=2mu\bm{z}=(\bm{R}\bm{W}^0\bm{D})^\top\bm{x}$\footnote{\textbf{Errata}: The NeurIPS camera ready version contained a typographical error in the forward pass for re-scaled OFT, listed as $\thickmuskip=2mu \medmuskip=2mu\bm{z}=(\bm{D}\bm{R}\bm{W}^0)^\top\bm{x}$. The originally implemented code is accurate, ensuring Appendix~\ref{app:rescaled} remains correct.}, where $\thickmuskip=2mu \medmuskip=2mu\bm{D}=\text{diag}(s_1,\cdots,s_n)\in\mathbb{R}^{n\times n}$ denotes a diagonal matrix with strictly positive, learnable entries $s_1,\ldots,s_n$. Unlike the OFT setup in Eq.~\eqref{eq:oft_general}, where the only trainable component is $\bm{R}$, this variant also treats $\bm{D}$ as a learnable parameter. This extension increases the expressive power of OFT with a negligible increase in parameter count. Although we use standard OFT for the main empirical validation to highlight the benefits of orthogonal transformations specifically, our findings indicate the re-scaled OFT generally achieves improved performance (see Appendix~\ref{app:rescaled}).

\textbf{No extra cost during inference}. In contrast to ControlNet, our OFT approach does not incur any additional computational burden during inference. During inference, it suffices to pre-multiply the trained orthogonal matrix $\bm{R}$ with the original weight matrix $\bm{W}^0$ to produce an updated weight matrix $\thickmuskip=2mu \medmuskip=2mu \bm{W}=\bm{R}\bm{W}^0$, resulting in inference performance identical to that of the unmodified model.

\textbf{Finetuning stability and convergence}. We begin by assessing the stability and speed of convergence during finetuning across DreamBooth, LoRA, OFT, and COFT, as summarized in Figure~\ref{fig:teaser} and Figure~\ref{fig:db_convergence}. Notably, both COFT and OFT maintain a high degree of stability throughout the finetuning of the diffusion model. Within just 400 iterations, DreamBooth and the OFT variants demonstrate robust control over generation, yet LoRA exhibits a deficiency in maintaining subject identity. At the 2000-iteration mark, DreamBooth starts producing degenerate outputs, and LoRA fails to accurately synthesize the yellow shirt, mistakenly producing yellow fur instead. On the other hand, OFT and COFT consistently offer reliable and stable control over the output images. These findings underscore the capacity of our OFT framework to deliver both rapid and stable convergence for subject-driven image synthesis. It is also important to highlight that insensitivity to finetuning iteration counts is crucial, as this feature reduces the need for tedious hyperparameter adjustment for different subjects. As such, both OFT and COFT allow for straightforward setting of a generous iteration budget without necessitating precise tuning. With COFT, given a suitable choice of $\epsilon$, parameters like learning rate and total iteration count become much easier to configure.


\begin{abstract}
State-of-the-art text-to-image diffusion models are capable of synthesizing highly realistic images in response to textual prompts. A pressing research question is how to steer these advanced models effectively to address a variety of downstream tasks. In this work, we present Orthogonal Finetuning (OFT), a methodical finetuning approach tailored for adapting text-to-image diffusion models to new tasks. Distinct from previous techniques, OFT is theoretically guaranteed to maintain hyperspherical energy—a property reflecting pairwise neuron interactions on the unit hypersphere. We observe that maintaining this property is essential to sustain the semantic fidelity of text-to-image diffusion models after adaptation. To further enhance the robustness of the finetuning process, we propose Constrained Orthogonal Finetuning (COFT), which incorporates a radius constraint on the hypersphere. Our study focuses on two key adaptation scenarios: (1) subject-driven generation, where given a limited set of subject images and a text description, the objective is to produce images featuring the same subject in varied contexts; (2) controllable generation, where the model leverages additional input control signals to guide the image synthesis process. Our empirical analysis demonstrates that the OFT approach yields superior image quality and more rapid convergence when compared to existing finetuning methods.
\end{abstract}

The convergence of the series in operator norm implies we can shift the constraint $\thickmuskip=2mu \medmuskip=2mu\|\bm{R}-\bm{I}\|\leq\epsilon$ into the Cayley transform, which translates to the new condition $\thickmuskip=2mu \medmuskip=2mu \|\bm{Q}-\bm{0}\|\leq\epsilon'$; here, $\bm{0}$ represents the zero matrix, and $\epsilon'$ is a distinct tuning parameter from $\epsilon$. This reformulated constraint on $\bm{Q}$ is straightforward to impose using projected gradient descent. For identity initialization of the orthogonal transformation $\bm{R}$, $\bm{Q}$ is simply set to be the zero matrix. Conceptually, COFT incorporates two direct constraints: minimizing hyperspherical energy variation and limiting deviation from the original pretrained parameters. Although many current finetuning methods employ this second constraint in an implicit way, COFT formalizes it as an explicit element. While OFT already achieves strong results, we find that COFT exhibits enhanced stability during finetuning, attributable to this explicit regularization. The relationship between the choice of $\epsilon$ and COFT’s performance is illustrated in Figure~\ref{fig:eps_coft}. As the value of $\epsilon$ increases, the behavior of COFT becomes more aligned with OFT, whereas smaller values of $\epsilon$ make COFT’s finetuned output closely resemble that of the pretrained text-to-image diffusion model.